% $Date: 2022/12/30 17:24:03 $
% This template file is public domain.
%
% TUGboat class documentation is at:
%   http://mirrors.ctan.org/macros/latex/contrib/tugboat/ltubguid.pdf
% or
%   texdoc tugboat

\documentclass{ltugboat}

\usepackage{microtype}
\usepackage{graphicx}
\usepackage[hidelinks,pdfa]{hyperref}

%%% Start of metadata %%%

\title{Example \TUB\ article}

% repeat info for each author; comment out items that don't apply.
\author{First Last}
\address{Street Address \\ Town, Postal \\ Country}
\netaddress{user (at) example dot org}
\personalURL{https://example.org/~user/}
%\ORCID{0}
% To receive a physical copy of the TUGboat issue, please include the
% mailing address we should use, as a comment if you prefer it not be printed.

%%% End of metadata %%%

\begin{document}
\maketitle

\begin{abstract}
Please write an abstract. Please use only standard \LaTeX\ and \TUB\
macros in the abstract, not article-specific abbreviations or other
macros; that helps us make the web pages. Please avoid bibliography
references in the abstract.
\end{abstract}

\section{Introduction}

This is an example article for \TUB, linked from
\url{https://tug.org/TUGboat/location.html}.

\section{Basic packages and hyperlinks}

The standard \texttt{graphicx} and \texttt{xcolor} packages work fine
for image and color handling, respectively, as does the
\texttt{hyperref} package for active urls. \TUB\ is produced using
\acro{PDF} files exclusively, but conversion from \DVI\ is fine and supported.

Please use \cs{tburl} for urls, as in:
\begin{verbatim}
\tburl{some/url}             % yes
\tburl{https://some/ftp/url} % yes
\tburl{ftp://some/ftp/url}   % yes
\end{verbatim}
For \TUB, we prefer not to print the `\texttt{https://}', but to include
the protocol in other cases. The \cs{tburl} macro takes care of all
cases (when \texttt{hyperref} is loaded). It's also fine to just use
\cs{url}.

In any case, please do \emph{not} hide hyperlinks via \cs{href}, as in:
\begin{verbatim}
\href{https://some/url}{some text} % no!
\end{verbatim}
When printed on physical paper, the url will not be visible; indeed,
readers will not know it exists.

\section{Abbreviation macros}

The \texttt{ltugboat} class provides many abbreviation commands; here
are a few of the most common:

% verbatim blocks are often better in \small
\begin{verbatim}[\small]
\AllTeX \AMS
\BibTeX \Cplusplus \CTAN
\DVI
\HTML \LaTeXe
\MacOSX \MathML \MF
\PDF \PS
\TUB \TUG \tug
\WEB
\XeLaTeX \XeTeX \XML

\Dash  % em-dash with our spacing
\slash % breakable /

\acro{FRED} -> {\small[er] FRED}  % please use
\cs{fred}   -> \fred
\meta{fred} -> <fred>
\titleref{Book of Fred}
\end{verbatim}

\section{Figures}

For \TUB, the standard \texttt{figure} environment produces a
column-width figure; this is desirable when at all possible. The
\texttt{figure*} environment produces a full-width (across both columns)
figure when needed. Analogously for \texttt{table} and \texttt{table*}.

Please put captions below figures, but above tables. That's our
convention.

Don't worry overmuch about figure placement, as this will likely change
with any edits.

\begin{figure}
A column-width figure, made with \cs{begin}\tubbraced{figure}.
\caption{Caption for column-width figure.}
\label{fig.colwidth}
\end{figure}

\begin{figure*}
A full-width figure, made with \cs{begin}\tubbraced{figure*}.
\caption{Caption for full-width figure.}
\label{fig.fullwidth}
\end{figure*}

\section{References}

For references to other issues of \TUB, please use the format
\textsl{volno}:\textsl{issno}, e.g., ``\TUB\ 32:1'' for volume~32,
number~1.

The primary \TUB\ style documentation is the \texttt{ltubguid} document
at \tbsurl{ctan.org/pkg/tugboat}. For general \CTAN\ package references,
we recommend the form just shown, using \texttt{/pkg/}; if you need to
refer to a specific file on \CTAN, use:\\
\texttt{https://mirror.ctan.org/\textsl{path}}

We recommend using \BibTeX\ (but don't require it), with the
\texttt{tugboat} \BibTeX\ style. The \texttt{biblio} directory on \CTAN,
\tbsurl{ctan.org/pkg/biblio}, provides files \texttt{tugboat.bib} with
an essentially complete bibliography of \TUB, and \texttt{texbook3.bib}
with many \TeX-related citations, along with many others. Both of these
(and many more) should also already be in a non-minimal \TeX\
distribution.

Email \verb|tugboat@tug.org| with any problems or questions.

\bibliographystyle{tugboat}
\nocite{book-minimal}      % make the bibliography non-empty
\bibliography{xampl}       % xampl.bib comes with BibTeX

\makesignature
\end{document}